\chapter{Implicación, iniciativa y toma de decisiones}

\section{Implicación profesional y ética}
El desarrollo del proyecto ha requerido un comportamiento profesional en varios sentidos. Se ha evitado publicar credenciales reales, se han versionado configuraciones saneadas, se han mantenido separados los servicios críticos y se ha documentado con honestidad tanto lo conseguido como lo que todavía no está cerrado. En lugar de maquillar limitaciones, se han tratado como parte de la evaluación técnica.

La dimensión ética es especialmente importante porque Odín opera sobre datos íntimos: fotos, notas, documentos, sensores y cámaras. Por ello se prioriza el control local, la reducción de exposición y la seguridad operativa frente a soluciones más vistosas pero menos responsables.

\section{Iniciativa}
El proyecto muestra iniciativa en varios frentes:
\begin{itemize}
  \item se ha pasado de una idea amplia a un sistema realmente desplegado;
  \item se han investigado múltiples motores de voz hasta comprender dónde estaba el cuello de botella real;
  \item se sustituyó la estrategia de acceso remoto cuando dejó de ser la mejor opción;
  \item se añadió monitorización visual y autoreparabilidad tras detectar necesidades operativas;
  \item se reorganizó una infraestructura que había crecido de forma experimental para convertirla en algo mantenible.
\end{itemize}

\section{Decisiones justificadas}
\begin{table}[H]
  \centering
  \caption{Ejemplos de decisiones justificadas}
  \label{tab:decisiones-justificadas}
  \begin{tabularx}{\textwidth}{lX}
    \toprule
    Decisión & Justificación \\
    \midrule
    Migrar a Tailscale & Menor carga operativa y mejor encaje con red doméstica y VPN privada. \\
    Mantener Home Assistant en Raspberry Pi 5 & Ya era una pieza estable y especializada; moverla a Docker no aportaba valor. \\
    Usar CPU en Frigate para detección & Rendimiento suficiente y menor competencia por la GPU AMD. \\
    Usar GPU en Immich y Ollama & Son cargas donde la aceleración aporta valor claro. \\
    Aceptar límite de voz & Evita consumir el proyecto en una optimización con retorno decreciente. \\
    Reorganizar Docker antes de seguir creciendo & Mejora mantenibilidad y prepara la documentación final. \\
    \bottomrule
  \end{tabularx}
\end{table}

\section{Respuesta ante obstáculos}
Los obstáculos no se han resuelto escondiéndolos, sino estudiándolos:
\begin{itemize}
  \item problemas de compatibilidad de TTS con AMD;
  \item fallos de Ollama investigados mediante logs;
  \item automatizaciones que requerían revisión;
  \item crecimiento desordenado de contenedores;
  \item necesidad de avisos útiles y no meramente cosméticos.
\end{itemize}

La actitud seguida ha sido coherente con un trabajo de ingeniería: diagnosticar, comparar, decidir, documentar y continuar. Esa secuencia es tan evaluable como el resultado técnico final.
